\chapter{Środowisko i metodyka}\label{ch:architektura_laboratorium}

\section{Topologia sieci i wirtualizacja}\label{sec:topologia_wirtualizacja}

Środowisko badawcze zostało zaprojektowane jako całkowicie odizolowana sieć wirtualna w środowisku \code{VMware Workstation}. Celem takiej konfiguracji jest uniemożliwienie wydostania się złośliwego kodu poza laboratorium przy jednoczesnym zapewnieniu próbkom odpowiedzi sieciowych niezbędnych do pełnego uruchomienia. Schemat połączeń w laboratorium przedstawiono na rysunku \ref{fig:topologia_lab}, natomiast konfigurację adresacji ujęto w tabeli \ref{tab:adresacja_sieciowa}.

\Figure[0.85]{topology.png}{Topologia połączeń w odizolowanej sieci laboratoryjnej}{fig:topologia_lab}

\begin{table}[H]
\centering
\caption{Konfiguracja adresacji sieci laboratoryjnej}
\label{tab:adresacja_sieciowa}
\tablefont
\begin{tabularx}{\linewidth}{l l L{3.8cm} Y}
\toprule
\hcell{Węzeł} & \hcell{Adres IP} & \hcell{Rola w sieci} & \hcell{Zadanie w środowisku} \\
\midrule
REMnux & \code{192.168.100.1} & Brama domyślna / DNS & Emulacja usług sieciowych, przechwyt ruchu i analiza zebranych artefaktów \\
Windows 10 x64 & \code{192.168.100.10} & Stacja ofiary & Uruchamianie badanych próbek i rejestracja telemetrii systemowej \\
\bottomrule
\end{tabularx}
\end{table}

Komunikacja pomiędzy maszynami odbywa się w wydzielonym segmencie sieci wirtualnej (\code{VMnetX}) z wyłączonym serwerem DHCP, zablokowanym translatorem adresów (NAT) oraz odpiętym wirtualnym interfejsem hosta. Dzięki temu stacja ofiary widzi w sieci lokalnej wyłącznie bramę \code{REMnux}.

Zastosowanie środowiska \code{VMware Workstation} zapewnia dwie korzyści istotne dla analizy śledczej:
\begin{itemize}
    \item \textbf{Bezagentowa akwizycja pamięci RAM:} wstrzymanie maszyny wirtualnej lub wykonanie migawki zapisuje na dysku hosta surowy obraz pamięci fizycznej gościa (\code{.vmem}). Pozwala to pobrać zawartość pamięci RAM z zewnątrz, bez uruchamiania dodatkowych programów na zainfekowanej stacji i bez modyfikowania stanu jej procesów.
    \item \textbf{Izolacja od hosta i maskowanie wirtualizacji:} na stacji ofiary zrezygnowano z instalacji pakietu dodatków \code{VMware Tools}, wyłączono współdzielenie schowka oraz folderów z systemem gospodarza, a w konfiguracji maszyny ograniczono widoczność podstawowych identyfikatorów wirtualizatora. Chroni to system hosta oraz zapobiega przedwczesnemu przerwaniu działania przez próbki wyposażone w proste mechanizmy wykrywania środowisk wirtualnych (\textit{Anti-VM}).
\end{itemize}

\section{Przygotowanie stacji Windows}\label{sec:konfiguracja_windows}

Na stacji ofiary z systemem \code{Windows 10 x64} wprowadzono zmiany konfiguracyjne mające na celu umożliwienie pełnego wykonania złośliwego oprogramowania oraz zarejestrowanie wszystkich jego działań w dziennikach systemowych.

W pierwszej kolejności wyłączono wbudowane mechanizmy ochronne, które mogłyby zablokować uruchomienie próbki lub usunąć jej komponenty w trakcie działania:
\begin{itemize}
    \item dezaktywowano ochronę w czasie rzeczywistym, monitorowanie behawioralne oraz ochronę przed naruszeniami w usłudze \code{Windows Defender},
    \item wyłączono blokady wykonywania skryptów \code{PowerShell} oraz skanowanie za pośrednictwem interfejsu \code{AMSI}, co pozwala na wykonanie obfuskowanych skryptów ładujących (\textit{loaders}),
    \item wyłączono zaporę systemową oraz usługę automatycznych aktualizacji, aby zapewnić swobodny ruch sieciowy w kierunku bramy \code{REMnux} i wyeliminować szum sieciowy generowany przez sam system operacyjny.
\end{itemize}

W drugim kroku skonfigurowano rozszerzoną rejestrację zdarzeń. Standardowe dzienniki systemu Windows nie rejestrują szczegółowych operacji na procesach i pamięci, dlatego na stacji zainstalowano usługę \code{Sysmon} \cite{russinovich_sysmon} z konfiguracją opartą na profilu \code{Sysmon-Modular} \cite{sysmon_modular}. Zestawienie najważniejszych monitorowanych kategorii zdarzeń przedstawiono w tabeli \ref{tab:sysmon_events}. Dodatkowo włączono rejestrowanie treści wykonywanych bloków skryptowych \code{PowerShell} oraz zwiększono maksymalny rozmiar buforów dzienników zdarzeń (\code{.evtx}), aby zapobiec nadpisaniu logów podczas intensywnej aktywności próbki.

\begin{table}[htbp]
\centering
\caption{Główne kategorie zdarzeń monitorowane przez usługę Sysmon}
\label{tab:sysmon_events}
\tablefont
\begin{tabularx}{\linewidth}{c L{4.2cm} Y}
\toprule
\hcell{Event ID} & \hcell{Kategoria zdarzenia} & \hcell{Zastosowanie w analizie poinfekcyjnej} \\
\midrule
1 & ProcessCreate & Odtworzenie drzewa procesów, argumentów wiersza poleceń i użycia narzędzi LOLBins. \\
3 & NetworkConnect & Powiązanie połączeń sieciowych TCP/UDP z konkretnym procesem (\code{PID}). \\
7 & ImageLoad & Wykrywanie ładowania nietypowych bibliotek DLL przez procesy. \\
8 & CreateRemoteThread & Identyfikacja tworzenia zdalnych wątków przy wstrzykiwaniu kodu. \\
10 & ProcessAccess & Wykrywanie otwierania pamięci obcych procesów (np. pod kątem iniekcji lub odczytu poświadczeń). \\
11 & FileCreate & Rejestracja plików wykonywalnych i skryptów zrzucanych na dysk. \\
12 / 13 & RegistryEvent & Śledzenie modyfikacji kluczy rejestru odpowiedzialnych za autostart. \\
22 & DNSEvent & Rejestracja zapytań o nazwy domenowe serwerów C2. \\
\bottomrule
\end{tabularx}
\end{table}

\section{Węzeł REMnux i narzędzia analityczne}\label{sec:konfiguracja_remnux}

Druga maszyna wirtualna, pracująca pod kontrolą dystrybucji \code{REMnux}, pełni jednocześnie funkcję lokalnego emulatora sieci oraz stacji roboczej do analizy zebranego materiału dowodowego.

Wiele próbek złośliwego oprogramowania przed uruchomieniem głównego ładunku weryfikuje dostępność połączenia sieciowego lub próbuje nawiązać sesję z serwerem dowodzenia (C2). Na węźle \code{REMnux} skonfigurowano usługę \code{INetSim}, która nasłuchuje na adresie bramy (\code{192.168.100.1}) i odpowiada na ruch przychodzący ze stacji ofiary:
\begin{itemize}
    \item przekierowuje wszystkie zapytania DNS na adres bramy analitycznej,
    \item obsługuje połączenia HTTP oraz HTTPS (po uprzednim dodaniu certyfikatu urzędu certyfikacji \code{INetSim} do magazynu zaufanych urzędów na stacji Windows, co zapobiega odrzucaniu połączeń szyfrowanych),
    \item przyjmuje połączenia na pozostałych standardowych portach usług aplikacyjnych (m.in. SMTP, FTP, surowe gniazda TCP).
\end{itemize}

Cały ruch przechodzący przez interfejs sieciowy węzła \code{REMnux} jest równolegle zapisywany do plików \code{.pcap} za pomocą narzędzia \code{tshark}. Zebrane zrzuty pamięci, dzienniki zdarzeń oraz pakiety sieciowe są następnie przetwarzane przy użyciu zestawu narzędzi podzielonego według zadań analitycznych (tabela \ref{tab:stos_narzedziowy}).

\begin{table}[htbp]
\centering
\caption{Zestaw narzędzi wykorzystywanych w analizie i inżynierii detekcji}
\label{tab:stos_narzedziowy}
\tablefont
\begin{tabularx}{\linewidth}{L{4.5cm} L{3.5cm} Y}
\toprule
\hcell{Obszar analizy} & \hcell{Narzędzia} & \hcell{Przeznaczenie} \\
\midrule
Analiza statyczna & PEstudio, FLOSS & Wstępna ocena struktury plików wykonywalnych PE, sekcji, tablicy importów oraz wyodrębnienie ciągów znaków. \\
\addlinespace
Analiza pamięci RAM & Volatility \cite{volatility} & Przegląd listy procesów, wykrywanie podejrzanych obszarów pamięci wykonywalnej, identyfikacja wstrzyknięć kodu oraz zrzucanie ładunków z RAM. \\
\addlinespace
Analiza logów zdarzeń & Hayabusa \cite{hayabusa_docs} & Szybkie przeszukiwanie plików \code{.evtx} za pomocą reguł detekcyjnych oraz generowanie chronologicznej osi czasu zdarzeń. \\
\addlinespace
Inżynieria detekcji (RAM i pliki) & YARA \cite{yara_docs} & Tworzenie i testowanie reguł dopasowujących wzorce bajtowe oraz tekstowe w zrzutach pamięci i wyodrębnionych plikach. \\
\addlinespace
Inżynieria detekcji (logi) & Sigma \cite{sigma_hq} & Definiowanie reguł wykrywających szkodliwe zachowania na podstawie zdarzeń z dzienników systemowych i usługi Sysmon. \\
\bottomrule
\end{tabularx}
\end{table}

\section{Przebieg badania i walidacja reguł}\label{sec:przebieg_badania}

Dla zachowania powtarzalności oraz spójności zebranych danych każda próbka złośliwego oprogramowania jest badana według tej samej, sześcioetapowej procedury przedstawionej w tabeli \ref{tab:cykl_badawczy}. Kolejność zabezpieczania artefaktów została ustalona zgodnie z zasadą kolejności ulotności (\textit{Order of Volatility}, RFC 3227 \cite{rfc3227}) - w pierwszej kolejności zamrażany jest stan maszyny wirtualnej i pobierany jest obraz pamięci RAM (\code{.vmem}), a dopiero w kolejnym kroku eksportowane są pliki dzienników z dysku (\code{.evtx}). Dzięki temu czynności związane z eksportem logów nie nadpisują zawartości pamięci operacyjnej.

\begin{table}[htbp]
\centering
\caption{Procedura przeprowadzenia eksperymentu i pobrania materiału dowodowego}
\label{tab:cykl_badawczy}
\tablefont
\begin{tabularx}{\linewidth}{c L{4.2cm} Y}
\toprule
\hcell{Etap} & \hcell{Nazwa etapu} & \hcell{Wykonywane czynności} \\
\midrule
1 & Przywrócenie stanu bazowego & Odtworzenie czystej migawki na stacji Windows oraz wyczyszczenie logów z poprzednich badań na węźle REMnux. \\
\addlinespace
2 & Uruchomienie nasłuchu & Włączenie symulacji usług w \code{INetSim} oraz rozpoczęcie zapisu ruchu sieciowego w \code{tshark}. \\
\addlinespace
3 & Uruchomienie próbki & Przeniesienie badanej próbki na stację ofiary, jej uruchomienie oraz odczekanie na wykonanie operacji w pamięci, systemie plików i sieci. \\
\addlinespace
4 & Zrzut pamięci RAM & Wstrzymanie maszyny wirtualnej lub utworzenie migawki z poziomu hiperwizora i zabezpieczenie pliku \code{.vmem} przed wykonaniem jakichkolwiek poleceń w systemie gościa. \\
\addlinespace
5 & Pobranie logów i ruchu & Zatrzymanie zapisu pliku \code{.pcap} na węźle REMnux oraz wyeksportowanie dzienników \code{.evtx} (Sysmon, Security, PowerShell) ze stacji Windows. \\
\addlinespace
6 & Obliczenie sum kontrolnych & Wyznaczenie skrótów SHA-256 dla zebranych plików dowodowych (\code{.vmem}, \code{.evtx}, \code{.pcap}) i przekazanie ich do analizy. \\
\bottomrule
\end{tabularx}
\end{table}

Po zebraniu materiału dowodowego wyniki z poszczególnych źródeł są ze sobą łączone (korelacja). Punktem wspólnym pomiędzy zrzutem pamięci RAM, logami \code{Sysmon} oraz ruchem sieciowym \code{.pcap} są identyfikatory procesów (\code{PID} oraz \code{ProcessGuid}), numery portów i adresy sieciowe oraz znaczniki czasu w strefie UTC. Pozwala to powiązać konkretny proces widoczny w logach systemowych z wstrzykniętym do jego pamięci kodem oraz ruchem sieciowym, który ten proces wygenerował.

Na podstawie zidentyfikowanych artefaktów tworzone są reguły detekcyjne \code{YARA} oraz \code{Sigma}. Każda reguła podlega dwuetapowej weryfikacji:
\begin{itemize}
    \item \textbf{Test wykrywalności (\textit{True Positive}):} sprawdzenie, czy reguła poprawnie wykrywa zagrożenie na zebranym materiale poinfekcyjnym (zrzucie pamięci lub logach \code{.evtx} z zainfekowanej stacji).
    \item \textbf{Test fałszywych alarmów (\textit{False Positive}):} przeskanowanie tą samą regułą zrzutu pamięci oraz dzienników zdarzeń pochodzących z czystego, niezainfekowanego systemu bazowego w celu potwierdzenia, że reguła nie generuje błędnych trafień na standardowych procesach systemowych.
\end{itemize}
